% Flächeninhalt durch Integration – bank/flaecheninhalt-durch-integration/ (zusammenbau v0.4, Heft)
\einheitenkopf[t3]{Flächeninhalt durch Integration}
\setcounter{aufgabe}{7}

% Anlauf (ohne Prüfkennung)
\begin{aufgabe}{Zusammensetzen, Maßstab, Volumen – Anlauf}
\begin{teile}
\teil Teilstück unter dem Bogen von $f(x) = 0{,}5x^2 + 1$ zwischen $x = 0$ und $x = 2$ – Integral oder Formelfläche? \kreuz{Integral}\\ \kreuz{Formelfläche}
\teil Eine Fläche wird oben vom Graphen von $f(x) = x^2$ ($0 \le x \le 2$) und der Strecke von $(2|4)$ nach $(4|0)$ begrenzt, unten von der x-Achse – Inhalt?
\teil Der obere Rand ist $f(x) = 0{,}5x + 1$ für $0 \le x \le 2$ und $g(x) = -0{,}5x^2 + 2x$ für $2 \le x \le 4$; unten die x-Achse, links die y-Achse – Inhalt?
\end{teile}
\end{aufgabe}

\begin{aufgabe}{Zusammensetzen, Maßstab, Volumen – Prüfungsaufgaben}
\begin{teile}
\teil $f(x) = -x^2 + x + 2$. Jemand rechnet: (1) $2x + 2 = f(x)$ liefert $-1$ und $0$; (2) das Integral von $-1$ bis $0$ über $f(x) - (2x + 2)$ ergibt $\frac{1}{6}$; (3) $0{,}5 \cdot f(0) = 1$; (4) die Zahlen werden ins Verhältnis gesetzt. Deute jeden Schritt geometrisch. \hfill (Abitur 2024 GK)
\teil $f(x) = -x^2 + 2x + 3$. Jemand rechnet: (1) $3x + 3 = f(x)$ liefert $-1$ und $0$; (2) das Integral von $-1$ bis $0$ über $f(x) - (3x + 3)$ ergibt $\frac{1}{6}$; (3) $0{,}5 \cdot f(0) = \frac{3}{2}$; (4) die Zahlen werden ins Verhältnis gesetzt. Deute jeden Schritt geometrisch. \hfill (Abitur 2024 GK)
\teil Eine Böschung folgt dem Graphen von $f(x) = 4 - x^2$ (eine Einheit entspricht $1$ m). Sie wird im Bereich $1 \le x \le 2{,}5$ bis zur Tangente an den Graphen in $x = 1$ aufgeschüttet, $6$ m breit. Wie viel Kubikmeter Erde sind nötig? \hfill (Abitur 2020 GK)
\teil Ein Wall folgt dem Graphen von $f(x) = 2 - 0{,}5x^2$ (eine Einheit entspricht $1$ m). Er wird im Bereich $1 \le x \le 2{,}5$ bis zur Tangente an den Graphen in $x = 1$ aufgeschüttet, $8$ m breit. Wie viel Kubikmeter Erde sind nötig? \hfill (Abitur 2020 GK)
\end{teile}
\end{aufgabe}

% Anlauf (ohne Prüfkennung)
\begin{aufgabe}{Flächenbedingungen – Anlauf}
\begin{teile}
\teil Gesucht ist der Inhalt der Fläche zwischen dem Graphen von $f(x) = x^2$ und der Geraden $y = 4$ – vorwärts oder rückwärts? \kreuz{vorwärts (Inhalt gesucht)}\\ \kreuz{rückwärts (Inhalt gegeben)}
\teil Für $m > 0$ schließt die Gerade $y = m \cdot x$ mit dem Graphen von $f(x) = 2x^2$ eine Fläche vom Inhalt $9$ ein – $m$?
\teil Die Fläche unter dem Graphen von $f(x) = 3x^2$ von $0$ bis $2$ soll von einer Parallelen zur y-Achse $x = k$ halbiert werden – $k$? Runde auf zwei Nachkommastellen.
\end{teile}
\end{aufgabe}

\begin{aufgabe}{Flächenbedingungen – Prüfungsaufgaben}
\begin{teile}
\teil Der Graph von $f$ in der Abbildung ist punktsymmetrisch zu $(0|3)$. Die Gleichung Integral von $k$ bis $k + 2$ über $f(x)$ $= 6$ hat genau eine Lösung mit $-1{,}5 \le k \le 0$. Gib sie ohne Rechnung an und begründe mit Eintragungen in der Abbildung. \hfill (Abitur 2025 GK)

\begin{ksys}[xmin=-3,xmax=3,ymin=-1,ymax=6,ablesen]\funktion{0.5*\x^3-2*\x+3}{f}\end{ksys}
\teil Der Graph von $f$ in der Abbildung ist punktsymmetrisch zu $(0|5)$. Die Gleichung Integral von $k$ bis $k + 4$ über $f(x)$ $= 20$ hat genau eine Lösung mit $-3 \le k \le 0$. Gib sie ohne Rechnung an und begründe mit Eintragungen in der Abbildung. \hfill (Abitur 2025 GK)

\begin{ksys}[xmin=-3,xmax=3,ymin=-1,ymax=9,ablesen]\funktion{-0.25*\x^3+\x+5}{f}\end{ksys}
\end{teile}
\end{aufgabe}

% Anlauf (ohne Prüfkennung)
\begin{aufgabe}{Flächenbilanz – Anlauf}
\begin{teile}
\teil Integral von $0$ bis $2$ über $f(x)$, Graph im Bild – Vorzeichen des Werts? Ist der Wert der Flächeninhalt? \kreuz{Wert positiv}\\ \kreuz{Wert negativ}\\ \kreuz{Wert null}

\begin{ksys}[xmin=-3,xmax=5,ymin=-4,ymax=4,xstep=1,ystep=1,ablesen]\funktion{\x^2-4}{f}\end{ksys}
\teil Integral von $0$ bis $4$ über $f(x)$ – bestimme den Wert durch Kästchenzählen am Bild.

\begin{ksys}[xmin=-1,xmax=5,ymin=-4,ymax=4,xstep=1,ystep=1,ablesen]\funktion{0.5*\x+1}{f}\end{ksys}
\teil Begründe am Bild, dass das Integral von $-2$ bis $2$ über $f(x)$ null ist.

\begin{ksys}[xmin=-3,xmax=5,ymin=-4,ymax=4,xstep=1,ystep=1,ablesen]\funktion{0.5*\x^3-2*\x}{f}\end{ksys}
\end{teile}
\end{aufgabe}

\begin{aufgabe}{Flächenbilanz – Prüfungsaufgaben}
\begin{teile}
\teil Der Graph von $f(x) = 0{,}25(x - 3)^3 + 2$ ist punktsymmetrisch zu seinem Wendepunkt $W(3|2)$. Die Gerade $g(x) = x - 1$ geht durch $W$ und schneidet den Graphen bei $x = 1$, $3$ und $5$. Begründe ohne Rechnung, dass das Integral von $1$ bis $5$ über $f(x)$ gleich $\frac{1}{2} \cdot (5 - 1) \cdot f(5)$ ist. \hfill (Abitur 2019 GK)
\teil Der Graph von $f(x) = \frac{1}{9}(x - 4)^3 + 3$ ist punktsymmetrisch zu seinem Wendepunkt $W(4|3)$. Die Gerade $g(x) = x - 1$ geht durch $W$ und schneidet den Graphen bei $x = 1$, $4$ und $7$. Begründe ohne Rechnung, dass das Integral von $1$ bis $7$ über $f(x)$ gleich $\frac{1}{2} \cdot (7 - 1) \cdot f(7)$ ist. \hfill (Abitur 2019 GK)
\end{teile}
\end{aufgabe}
